% =============================================================================
% Assignments/CS502/08-cpu-scheduling-basics-assignment.tex
% Assignment 8: CPU Scheduling --- Concepts, Criteria, FCFS and SJF
%
% Student set — NO answers. Solution key: 08-cpu-scheduling-basics-solutions.tex (never published).
% Sourced from Slides/CS502/08-cpu-scheduling-basics.tex; every problem reuses
% the deck's notation (AT, tau, CT, TAT, WT, RT; hat-tau for a prediction) and
% extends its worked examples to new workloads. Every schedule and number was
% verified with a simulator (FCFS, non-preemptive SJF, SRTF; no ties occur).
%
% Compile with: cd Assignments/CS502 &&
%   TEXINPUTS=../../Preambles:$TEXINPUTS xelatex -interaction=nonstopmode 08-cpu-scheduling-basics-assignment.tex
%   BIBINPUTS=../..:$BIBINPUTS bibtex 08-cpu-scheduling-basics-assignment
%   TEXINPUTS=../../Preambles:$TEXINPUTS xelatex -interaction=nonstopmode 08-cpu-scheduling-basics-assignment.tex (x2)
%   (Windows/MiKTeX uses ; not : as the path separator)
% =============================================================================

\documentclass[11pt]{article}
\usepackage[margin=1in]{geometry}
\input{../../Preambles/header}
\coursecode{CS502}
\renewcommand{\thesection}{8.\arabic{section}}

\setlength{\emergencystretch}{2.5em}

\title{Assignment 8: CPU Scheduling --- Concepts, Criteria, FCFS and SJF}
\author{Auqib Hamid Lone\\[0.3em]
  \emph{Department of Computer Science \& Engineering}, \emph{GCET Ganderbal}}
\date{\today}

\begin{document}
\maketitle

\noindent\textbf{Deadline:} \key{two weeks from the date of issue} --- there is no automatic late penalty; if you need more time, ask before the deadline and an extension will be granted (see the course policy in the syllabus).

\noindent Answer every question. Show all working for Numerical and Design questions --- a bare final answer without the Gantt chart and the derivation earns partial credit only. All times are in milliseconds. Use $\mathrm{TAT} = \mathrm{CT} - \mathrm{AT}$, $\mathrm{WT} = \mathrm{TAT} - \tau$, $\mathrm{RT} = \mathrm{ST} - \mathrm{AT}$.

\noindent\textbf{Reading:} Silberschatz et al., \emph{Operating System Concepts}, Ch.~5 \S5.1--5.3.2, pp.~183--192 \cite{SilberschatzOSC_operating_system_concepts}; Stallings, \emph{Operating Systems: Internals and Design Principles}, Ch.~9 \S9.1--9.2, pp.~406ff.\ \cite{StallingsOS_internals_design_principles}.

\section{Conceptual}

\textbf{Q1.} Scheduling decisions, preemption, and the dispatcher \cite{SilberschatzOSC_operating_system_concepts}.
\begin{enumerate}
  \item[(a)] List the four circumstances under which a CPU-scheduling decision may be needed. For each, say whether a non-preemptive scheduler must make a decision there.
  \item[(b)] A process is running when (i) it calls \texttt{wait} for a child, (ii) a timer interrupt fires, (iii) a disk read issued earlier by another process completes, (iv) it calls \texttt{exit}. For each event, name the state transition and state whether FCFS, non-preemptive SJF, and SRTF may change the running process at that moment.
  \item[(c)] Name the three things the dispatcher does, and explain why a policy that switches more often lowers CPU utilization even when every process is ready all the time.
  \item[(d)] Explain why preemption creates a problem that non-preemptive scheduling avoids, with reference to shared data.
\end{enumerate}

\textbf{Q2.} FCFS, SJF, and their failure modes.
\begin{enumerate}
  \item[(a)] Describe the convoy effect in terms of one CPU-bound and several I/O-bound processes. Explain why both the CPU and the I/O devices end up under-utilized.
  \item[(b)] For $n$ processes all ready at time 0 run in the order $\tau_{(1)}, \ldots, \tau_{(n)}$, derive $\sum_i \mathrm{WT}_i = \sum_{k=1}^{n}(n-k)\tau_{(k)}$. Then show that swapping an adjacent long--short pair ($L$ then $S$, $L > S$) into short--long order lowers the sum by exactly $L - S$, and conclude that SJF minimizes average WT.
  \item[(c)] Explain why SJF cannot be implemented exactly at the level of short-term scheduling, and what the scheduler uses instead.
  \item[(d)] Explain starvation under SJF and SRTF, and why FCFS cannot starve a process.
\end{enumerate}

\section{Numerical}

\textbf{Q3.} Five processes have (AT, $\tau$) in ms: \texttt{P1} $(0,4)$, \texttt{P2} $(1,2)$, \texttt{P3} $(3,1)$, \texttt{P4} $(10,2)$, \texttt{P5} $(11,5)$. Ties (if any) are broken by earlier arrival, and SRTF preempts only for a strictly shorter remaining time.
\begin{enumerate}
  \item[(a)] Draw the Gantt chart under FCFS, non-preemptive SJF, and SRTF. Mark idle intervals and arrivals.
  \item[(b)] For each policy, tabulate CT, TAT, WT, and RT per process and give the averages.
  \item[(c)] Compute the CPU utilization over $[0, \text{last completion}]$, ignoring dispatch latency. Why is it the same for all three policies?
  \item[(d)] Count the dispatches made by each policy. With a dispatch latency of 0.5~ms per dispatch, how much CPU time does SRTF spend dispatching?
\end{enumerate}

\textbf{Q4.} Burst prediction and dispatch cost.
\begin{enumerate}
  \item[(a)] A process has $\hat{\tau}_0 = 6$~ms and $\alpha = 0.75$. Its first three measured CPU bursts are 10, 2, and 8~ms. Compute $\hat{\tau}_1$, $\hat{\tau}_2$, $\hat{\tau}_3$, showing each step.
  \item[(b)] Repeat (a) with $\alpha = 0$ and with $\alpha = 1$, and state in one sentence what each extreme means.
  \item[(c)] Expand $\hat{\tau}_3$ in terms of $\tau_0, \tau_1, \tau_2, \hat{\tau}_0$ for general $\alpha$ and verify that the weights sum to 1.
\end{enumerate}

\section{Design}

\textbf{Q5.} The lab server receives five jobs, (AT, $\tau$) in ms: \texttt{P1} $(0,9)$, \texttt{P2} $(1,4)$, \texttt{P3} $(2,2)$, \texttt{P4} $(4,1)$, \texttt{P5} $(6,3)$. You must recommend a policy for this server.
\begin{enumerate}
  \item[(a)] Produce the FCFS, non-preemptive SJF, and SRTF schedules (Gantt charts and per-process tables).
  \item[(b)] Build a comparison table of average WT, average TAT, average RT, and number of dispatches.
  \item[(c)] With a dispatch latency of 0.5~ms, compute the CPU utilization of FCFS and of SRTF over the whole run.
  \item[(d)] Identify which process is hurt most by SRTF relative to FCFS, explain why, and describe an arrival stream that would starve it.
  \item[(e)] Recommend a policy for students who mostly submit short compile jobs but occasionally a long simulation. Justify your choice against the criteria of the lecture.
\end{enumerate}

\textbf{Q6.} SJF is only as good as its prediction.
\begin{enumerate}
  \item[(a)] Two processes are ready at time 0. The exponential average predicts 9~ms for \texttt{X} and 5~ms for \texttt{Y}; their actual next bursts are 3~ms and 7~ms. Give the order SJF chooses, the resulting average WT, and the average WT of the ideal order.
  \item[(b)] Design a three-process workload, all arriving at different times, for which FCFS and non-preemptive SJF produce identical schedules. State the general condition that makes this happen.
  \item[(c)] For the workload of Q5, design one additional process (give its AT and $\tau$) whose arrival makes SRTF preempt \texttt{P5}. Justify with the remaining times at the arrival instant.
\end{enumerate}

\bibliography{../../Bibliography_base}
\bibliographystyle{plain}
\end{document}
